% ---------------------------------------------------------------------------
% Figure -- where the three interventions apply in the VLA inference pipeline.
% Added 2026-08-26 to give Section III visual weight and to make rho, the
% eligibility window, and the horizon h=cr legible at a glance, rather than
% padding the Methods prose.
%
% Drawn in TikZ, NOT matplotlib, on purpose: a vector schema has no Type 3
% fonts, so it passes IEEE PDF eXpress, the check that would have failed the
% old matplotlib heat map. Nothing to regenerate; edit here.
%
% Preamble needs:
%   \usepackage{tikz}
%   \usetikzlibrary{positioning, arrows.meta, fit, backgrounds}
%
% ⚠️ COMPILE-CHECK in Overleaf before trusting the layout: this was authored
% without a local LaTeX run. The structure is deliberately simple (one row of
% boxes, one feedback arrow above, three tags below) to keep it robust.
%
% Spans both columns (figure*), placed near the top of Section III.
% ---------------------------------------------------------------------------

\begin{figure*}[t]
\centering
\begin{tikzpicture}[
    font=\small,
    >=Stealth,
    box/.style={draw, rounded corners=2pt, minimum height=9mm,
                minimum width=17mm, align=center, inner sep=3pt},
    tag/.style={draw, dashed, rounded corners=2pt, fill=black!5,
                align=center, font=\footnotesize, inner sep=3pt},
    flow/.style={->, thick},
]
% --- main pipeline, left to right ---
\node[box] (obs)  {Observation\\$o_t$};
\node[box, right=9mm of obs]  (enc)  {Visual\\encoder};
\node[box, right=9mm of enc]  (dec)  {Decoder\\$N$ layers};
\node[box, right=9mm of dec]  (head) {Action\\head};
\node[box, right=13mm of head] (env) {Environment};

\draw[flow] (obs)  -- (enc);
\draw[flow] (enc)  -- (dec);
\draw[flow] (dec)  -- (head);
\draw[flow] (head) -- node[above, font=\footnotesize]
            {chunk $(a_t^1,\dots,a_t^c)$} (env);

% --- feedback arrow above, carrying the horizon ---
\draw[flow] (env.north) -- ++(0,6mm)
            -| node[above, pos=0.25, font=\footnotesize]
            {next query after $h = c\,r$ steps} (obs.north);

% --- the three interventions, tagged at the point each acts ---
\node[tag, below=8mm of obs]  (fov) {\textbf{Visual foveation}\\warp before the\\encoder, keep $\rho$};
\node[tag, below=8mm of dec]  (dp)  {\textbf{Depth pruning}\\remove $k$ of $N$\\in a deep window};
\node[tag, below=8mm of head] (ar)  {\textbf{Action repetition}\\hold each action\\$r$ steps};

\draw[->, dashed] (fov) -- (obs);
\draw[->, dashed] (dp)  -- (dec);
\draw[->, dashed] (ar)  -- (head);
\end{tikzpicture}
\caption{The three interventions in the VLA inference pipeline. Each acts at
one point and moves one quantity: foveation transforms the observation before
the encoder, depth pruning removes layers inside the decoder, and action
repetition sets how many environment steps one query covers. The horizon
$h = c\,r$ combines the checkpoint's chunk length $c$ with the action repeat
$r$, so a single value of $r$ is a different horizon on each backbone.}
\label{fig:pipeline}
\end{figure*}
